\documentclass[10pt,a4paper]{amsart}
\usepackage[margin=19mm]{geometry}
\usepackage{amsmath,amssymb,mathtools}
\usepackage[scr=rsfs,cal=euler]{mathalfa}
\usepackage{xfrac}
\usepackage[unicode,hidelinks,bookmarks=false]{hyperref}
\newcommand{\ie}{i.e.\ }
\newcommand{\cf}{cf.\ }
\newcommand{\resp}{respectively\ }
\newcommand{\Subsection}{\subsection}
\providecommand{\nobreakdash}{\nobreak\hbox{-}}
\setlength{\emergencystretch}{2em}
\multlinegap0pt
\usepackage{mathtools}
\usepackage{tikz}
\usepackage{xcolor}
\newtheorem{theorem}{Theorem}
\newtheorem{corollary}{Corollary}
\newtheorem{lemma}{Lemma}
\newcommand{\B}{\mathcal{B}}
\renewcommand{\H}{\mathcal{H}}
\title{Operator Geometry of Hilbert Ball Automorphisms}
\author{Saikat Roy}
\date{Source: arXiv:2606.18736v1 (CC BY 4.0)}
\begin{document}
\maketitle

\noindent\emph{Source and licence.} Operator Geometry of Hilbert Ball Automorphisms. Version arXiv:2606.18736v1 (CC BY 4.0), distributed by arXiv under the Creative Commons Attribution 4.0 International licence, http://creativecommons.org/licenses/by/4.0/, as recorded in the arXiv licence metadata for this exact version and shown on the abstract page https://arxiv.org/abs/2606.18736v1. Full licence notice: \texttt{licenses/arxiv-2606.18736-CC-BY-4.0.txt}.\par\smallskip

\noindent\textit{Editorial scope.} Counted unit: the article's theorem Aut (source0.tex lines 138--140). The article's complete original proof of it is delivered with it (source0.tex lines 340--459, one \texttt{proof} environment of 120 lines, which the article places away from the statement). No mathematics is written, repaired or re-derived: the statement and the proof are the article's own lines, in the article's own notation, and no retained line is substituted or re-flowed.  Retained uncounted as the source-local context the proof needs: lines 222--231; lines 234--241.  Nothing was omitted from any retained span, so the package carries no omission record.  No retained line contains a citation, so the wrapper carries no bibliography and no bibliography entry.  No retained line contains a figure environment, an image-inclusion command or drawing code, so no image is delivered and the package declares no asset dependency.  Editorial formatting only, in addition: an AMS \texttt{amsart} preamble that restates the article's own declarations and macros used by the retained lines, namely the environments \texttt{theorem}, \texttt{corollary}, \texttt{lemma} on separate counters, together with the standard packages those lines need.  The retained environments print in this document as Theorem 1, Corollary 1, Lemma 1 in order of appearance, whereas the article prints the same items with the numbering of its own full document.  The complete native source of the article as distributed in the arXiv e-print for this exact version is bundled unchanged as \texttt{sources/203096/source0.tex}.
\begin{corollary}\label{smoothness-Hilbert}
Let $A$ be a nonzero operator acting on a Hilbert space $\H$. Then $A$ is a smooth point of 
$\mathcal{B}(\mathcal{H})$ if and only if the collection
\[
\left\{\mu \in \mathbb{C} : \lim_n \langle Tx_n, y_n \rangle 
= \mu,\ (x_n), (y_n) \subset S_{\mathcal{H}},\ 
\lim_n \langle Ax_n, y_n \rangle = \|A\|\right\}
\]
is a singleton set for every $T \in \mathcal{B}(\mathcal{H})$.
\end{corollary}

\begin{lemma}\label{block matrix norm}
Let $\mathcal{H}_0$ be a closed proper subspace of $\mathcal{H}$, and let $A \in \B(\mathcal{H})$ be such that $\mathcal{H}_0$ is a reducing subspace of $A$. Then 
\[
\|A\| = \max\left\{\|AP_0\|, \|A(I-P_0)\|\right\},
\]
where $P_0$ denotes the orthogonal projection onto $\mathcal{H}_0$. Moreover, if $\|A\| = \|AP_0\| > \|A(I-P_0)\|$, then $(z_n)$ is a norming sequence for $A$ if and only if $(z_n)$ is a norming sequence 
for $AP_0$.
\end{lemma}

\begin{theorem}\label{smooth: Aut}
Every element of $M \left(\H \right)$ which is not a pure rotation, represents a smooth point in $B \left(\H\oplus \mathbb{C} \right)$. 
\end{theorem}

\begin{proof}[Proof of Theorem \ref{smooth: Aut}]
Consider $M_{\phi_{a,U}} \in M(\mathcal{H})$ for some nonzero 
$a \in B$ and unitary $U$ on $\mathcal{H}$. We write $M_{\phi_{a,U}} = \widetilde{U}M_{a,I},$, where $\widetilde{U} = \mathrm{diag}(U, 1) \in \mathcal{B}(\mathcal{H} \oplus \mathbb{C})$ is unitary. Since $\widetilde{U}$ is unitary, $\|M_{\phi_{a,U}}\| = \|M_{a,I}\|$, and 
$(z_n)$ is a norming sequence for $M_{\phi_{a,U}}$ if and only if it is 
a norming sequence for $M_{a,I}$. 

\medskip

\noindent It is not difficult to see that $\mathcal{H}_0 = \left\{(z, 0) \in \mathcal{H} \oplus \mathbb{C} : z \perp a\right\}$ is a subspace of co-dimension $2$ in $\mathcal{H} \oplus \mathbb{C}$, and is an eigenspace of $M_{a,I}$ corresponding the eigenvalue $\alpha(a)$. Since $M_{a,I}$ 
is self-adjoint, and $\H_0$ is an invariant subspace of $\H$, $\mathcal{H}_0$ is in fact 
a reducing subspace of $M_{a,I}$, giving the block decomposition $M_{a,I} = T\oplus S$ with respect to the orthogonal decomposition $\mathcal{H} \oplus \mathbb{C} = \mathcal{H}_0 \oplus \mathcal{H}_0^\perp$, where $T = \alpha(a)I$ 
on $\mathcal{H}_0$ and $S = M_{a,I}|_{\mathcal{H}_0^\perp}$. 


\medskip

\noindent The complementary subspace $\mathcal{H}_0^\perp$ is 
a two-dimensional subspace, and spanned by the orthonormal basis 
$\left\{\frac{1}{\|a\|}(a, 0),\ (0,1)\right\}$. A direct computation shows that the matrix of $S$ with respect to this basis is given by
\[
[S] = \begin{bmatrix} 1 & -\|a\| \\ -\|a\| & 1 \end{bmatrix}.
\]
Since $[S]$ is self-adjoint, $\|S\|$ equals to the absolute value of its its largest eigenvalue, which is $1 + \|a\|$. The corresponding unit
eigenvector is 
\[
\frac{1}{\sqrt{2}}\left(\frac{1}{\|a\|}(a,0) - 
(0,1)\right) = \frac{1}{\sqrt{2}}\left(\frac{a}{\|a\|}, -1\right).
\]

\medskip

\noindent By Lemma~\ref{block matrix norm}, since 
$\|S\| = 1 + \|a\| > \alpha(a) = \|T\|$, we conclude
\[
\|M_{a,I}\| = \|S\| = 1 + \|a\|,
\]
and the norm is attained precisely on the unit sphere of the one-dimensional eigenspace:
\[
\left\{\mu\cdot\frac{1}{\sqrt{2}}
\left(\frac{a}{\|a\|}, -1\right) : |\mu| = 1\right\}.
\]
Moreover, by Lemma~\ref{block matrix norm} $(z_n)$ is a 
norming sequence for $M_{a,I}$ if and only if it is a norming 
sequence for $M_{a,I}P$, where $P$ denotes the orthogonal projection 
onto $\mathcal{H}_0^\perp$.



\medskip

Next, for any $A \in \mathcal{B}(\mathcal{H} \oplus \mathbb{C})$, consider the 
collection 
\begin{multline*}
\Lambda(M_{\phi_{a,U}}, A) = \left\{\mu \in \mathbb{C} : 
\lim_n \langle Az_n, y_n\rangle = \mu,\ 
(z_n),(y_n) \subset S_{\mathcal{H}\oplus\mathbb{C}},\right.\\
\left.\hspace{4cm} \lim_n \langle M_{\phi_{a,U}}z_n, y_n\rangle 
= \|M_{\phi_{a,U}}\|\right\}.
\end{multline*}
We show this collection is a singleton set. Let $\lambda_0 \in 
\Lambda(M_{\phi_{a,U}}, A)$, realized by the sequences $(z_n)$ and $(y_n)$. 
Since $|\langle M_{\phi_{a,U}}z_n, y_n\rangle| \leq \|M_{\phi_{a,U}}z_n\|$, the 
condition $\lim_n\langle M_{\phi_{a,U}}z_n, y_n\rangle = \|M_{\phi_{a,U}}\|$ 
forces $(z_n)$ to be a norming sequence for $M_{\phi_{a,U}}$, hence for 
$M_{a,I}$, and therefore, for $M_{a,I}P$.

\medskip

\noindent Since $\mathcal{H}_0^\perp$ is 
two-dimensional, its closed unit ball is a compact subset of $B_{\H\oplus \mathbb{C}}$. We observe that 
$\lim_n\|(I-P)z_n\| = 0$ (by the same argument as in 
Lemma~\ref{block matrix norm}), and the sequence $(Pz_n)$ lies in the closed unit ball of $\mathcal{H}_0^\perp$. We may therefore pass to a 
subsequence and assume $Pz_{n_k} \to z_0$ for some unit vector $z_0$ in $\mathcal{H}_0^\perp$. Since $(I-P)z_{n_k} \to 0$, we get 
$z_{n_k} \to z_0$, and we know from the norm attainment set of $M_{a,I}$ that the unit vector $z_0$ is uniquely determined by:
\[
z_0 = \mu_0\cdot\frac{1}{\sqrt{2}}\left(\frac{a}{\|a\|}, -1\right)
\quad\text{for some unimodular scalar } \mu_0.
\]

\noindent A parallel argument applies to $(y_n)$. Since $M_{a,I}$ 
is self-adjoint and $\lim_n\langle M_{\phi_{a,U}}z_n, y_n\rangle = 
\|M_{\phi_{a,U}}\| \in \mathbb{R}$, we have
\[
\lim_n\langle M_{\phi_{a,U}}z_n, y_n\rangle 
= \lim_n\langle M_{a,I}z_n, \widetilde{U}^*y_n\rangle 
= \lim_n\langle M_{a,I}\widetilde{U}^*y_n, z_n\rangle 
= \|M_{\phi_{a,U}}\|,
\]
which shows $(\widetilde{U}^*y_n)$ is also a norming sequence for 
$M_{a,I}$. By the same compactness argument, passing to a further 
subsequence gives $\widetilde{U}^*y_{n_k} \to y_0$ where
\[
y_0 = \sigma_0\cdot\frac{1}{\sqrt{2}}\left(\frac{a}{\|a\|}, -1\right)
\quad\text{for some unimodular scalar } \sigma_0.
\]

\noindent By a diagonal argument we extract a common subsequence 
$(n_p)$ along which both $z_{n_p} \to \mu_0\tilde{z}_0$ and 
$\widetilde{U}^*y_{n_p} \to \sigma_0\tilde{z}_0$, where $\mu_0$ and $\sigma_0$ are unimodular scalars and $\tilde{z}_0 = \frac{1}{\sqrt{2}}\left(a/\|a\|, -1\right)$. 
The condition $\lim_p\langle M_{\phi_{a,U}}z_{n_p}, y_{n_p}\rangle = 
\|M_{\phi_{a,U}}\|$ then gives
\[
\mu_0\overline{\sigma_0}(1 + \|a\|) = \|M_{\phi_{a,U}}\|,
\]
which forces $\mu_0\overline{\sigma_0} = 1$, hence $\mu_0 = \sigma_0$.

\noindent Therefore $\lambda_0$ is completely determined:
\[
\lambda_0 = \lim_p\langle Az_{n_p}, y_{n_p}\rangle 
= \left\langle A\tilde{z}_0,\ \widetilde{U}\tilde{z}_0\right\rangle
= \frac{1}{2}\left\langle A\!\left(\tfrac{a}{\|a\|}, -1\right),\ 
\widetilde{U}\!\left(\tfrac{a}{\|a\|}, -1\right)\right\rangle.
\]
Since this value depends only on $A$ and not on the choice of 
sequences, $\Lambda(M_{\phi_{a,U}}, A)$ is a singleton set for every 
$A \in \mathcal{B}(\mathcal{H} \oplus \mathbb{C})$. By 
Corollary~\ref{smoothness-Hilbert}, $M_{\phi_{a,U}}$ is a smooth point 
in $\mathcal{B}(\mathcal{H} \oplus \mathbb{C})$, completing the 
proof.
\end{proof}

\end{document}
